% One-page resume -- general software engineering roles.
%
% The research-oriented CV (AI / formal methods) is cv.tex.
% The full academic CV is academic-cv.tex.

\documentclass[letterpaper,10pt]{article}

\usepackage{geometry}
\usepackage{xcolor}
\usepackage[T1]{fontenc}
\usepackage[sc,osf]{mathpazo}

\def\name{Shaun Allison}

% Unfilled items render in red so they cannot ship unnoticed.
\newcommand{\todo}[1]{\textcolor{red}{\textbf{[TODO: #1]}}}

\geometry{
  letterpaper,
  top=0.6in, bottom=0.6in, left=0.75in, right=0.75in
}

\pagestyle{empty}

\usepackage{sectsty}
\sectionfont{\rmfamily\mdseries\large\scshape}

\setlength\parindent{0em}

\usepackage{enumitem}
\setlist[itemize,1]{label={--}, leftmargin=1.1em, topsep=0.15em, itemsep=0.1em,
                    parsep=0em}

% Tighten section spacing to hold one page
\usepackage{titlesec}
\titlespacing*{\section}{0pt}{6pt}{3pt}
\titleformat{\section}{\rmfamily\mdseries\large\scshape}{}{0em}{}[\vspace{-4pt}\rule{\linewidth}{0.4pt}]

% Job heading: title, org, dates
\newcommand{\role}[3]{\textbf{#1} \textnormal{---} #2 \hfill #3\par\vspace{1pt}}

\usepackage{hyperref}
\hypersetup{
  colorlinks=true,
  urlcolor=[rgb]{0.10,0.20,0.50},
  pdfauthor={\name},
  pdftitle={\name: Resume},
  pdfsubject={Resume}
}

\begin{document}

\begin{center}
  {\LARGE \name}\\[3pt]
  shallison@protonmail.com \quad\textbullet\quad
  \url{https://shalliso.github.io/} \quad\textbullet\quad
  \url{https://github.com/shalliso}\\[1pt]
  Toronto, Canada \quad\textbullet\quad US citizen and Canadian permanent
  resident
\end{center}

\vspace{2pt}

Software engineer with a computer science degree from Carnegie Mellon and a
PhD in mathematics. Built machine-learned ad and notification targeting at
Meta, holding two granted patents from that work, then spent a decade in
mathematical research --- most recently writing large Lean 4 formalizations and
publishing in top mathematics journals. Returning to industry engineering.

\section{Skills}

\begin{itemize}
  \item \textbf{Languages:} Python, Lean 4 (current); Hack, JavaScript,
        Haskell, C (industry use, 2012--2016)
  \item \textbf{Tools:} git, Mathlib, \LaTeX{}; familiar with numpy, pandas,
        matplotlib, PyTorch
  \item \textbf{Domains:} formal verification, machine-learned ranking and
        targeting, mathematical logic
\end{itemize}

\section{Experience}

\role{Software Engineer}{Meta (then Facebook)}{June 2015 -- August 2016}
\begin{itemize}
  \item Built machine-learned targeting for ads and notifications on the Games
        team --- deciding which users to reach and with what.
  \item Work produced two granted patents on machine-learned ranking (below).
\end{itemize}

\vspace{3pt}
\role{Postdoctoral Fellow / Lecturer}{Univ.\ of Toronto; Hebrew Univ.\ of
Jerusalem}{2021 -- 2025}
\begin{itemize}
  \item Independent research program in mathematical logic; published in the
        \emph{Journal of the European Mathematical Society}, \emph{Advances in
        Mathematics}, and the \emph{Journal of Mathematical Logic}.
  \item Instructor of record for eight courses, including a 400-student
        calculus class --- scoping a syllabus, building assessments, and
        managing a teaching staff.
\end{itemize}

\vspace{3pt}
\textbf{Independent Research} \hfill January 2026 -- present\par\vspace{1pt}
\begin{itemize}
  \item Two Lean 4 formalization projects (below); carried two journal papers
        through revision to acceptance.
  \item Writing a textbook on mathematical logic with Spencer Unger, alongside
        continuing research in descriptive set theory.
  \item Research visits to Cornell University and the University of Vienna;
        upcoming visit to the Institute for Advanced Study in Mathematics,
        Zhejiang University.
\end{itemize}

\vspace{4pt}
\emph{Earlier:} software engineering internship at Meta (2014); research and
development at VMware (2013); Cisco (2012).

\section{Projects}

\textbf{Knight model formalization} --- Lean 4, Mathlib, 2026 \hfill
\url{https://github.com/shalliso/KnightModel}
\begin{itemize}
  \item $\sim$5,700 lines of Lean formalizing a generalization of a classical
        model-theoretic construction. Developed with substantial AI assistance
        (Claude) under my direction; the mathematics and proof architecture are
        my own.
\end{itemize}

\vspace{3pt}
\textbf{Automatic continuity} --- Lean 4, Mathlib, 2025 \hfill
\url{https://github.com/shalliso/automatic_continuity}
\begin{itemize}
  \item $\sim$300 lines of Lean, written by hand. Complete machine-checked
        proof of a classical theorem in topological group theory.
\end{itemize}

\section{Patents}

\begin{itemize}
  \item \emph{Automatic recipient targeting for notifications} (US10506055B2)
        --- machine-learned selection of notification recipients. Granted.
  \item \emph{Selecting assets} (US11019177B2) --- machine-learned selection of
        which creative asset to serve. Granted.
\end{itemize}

\section{Education}

\begin{itemize}
  \item \textbf{Ph.D. Mathematics}, Carnegie Mellon University, 2021
  \item \textbf{B.S. Computer Science} and \textbf{B.S. Mathematics}
        (Discrete Mathematics and Logic), Carnegie Mellon University, 2015
  \item Putnam Competition top 500; CMU College Honors (Computer Science);
        Phi Beta Kappa
\end{itemize}

\end{document}
